% Esqueleto do artigo -- template SBC
% Precisa dos arquivos sbc-template.sty e sbc.bst (vêm no template oficial da SBC;
% no latex.sistemas.udesc.br, crie o projeto a partir do template SBC ou envie esses arquivos).
% Envie também: referencias.bib, tabela_descritiva.tex, boxplot_cc.pdf, boxplot_loc.pdf
\documentclass[12pt]{article}
\usepackage{sbc-template}
\usepackage{graphicx,url}
\usepackage[utf8]{inputenc}
\usepackage[T1]{fontenc}
\usepackage[brazil]{babel}
\usepackage{amsmath}

\sloppy

\title{Complexidade Ciclomática de Código Gerado por Modelos de Linguagem:\\
Um Experimento Comparativo entre [IA A] e [IA B]}

\author{Nome Aluno 1\inst{1}, Nome Aluno 2\inst{1}}

\address{Departamento de Engenharia de Software -- Centro de Educação Superior do Alto Vale do Itajaí\\
Universidade do Estado de Santa Catarina (UDESC) -- Ibirama, SC -- Brasil
  \email{\{aluno1,aluno2\}@edu.udesc.br}
}

\begin{document}
\maketitle

\begin{abstract}
% Escrever por último, em inglês, até 10 linhas: contexto, objetivo, método, resultado principal, conclusão.
[TO WRITE]
\end{abstract}

\begin{resumo}
% Mesmo conteúdo do abstract, em português.
[ESCREVER POR ÚLTIMO]
\end{resumo}

%=====================================================================
\section{Introdução}
% Parágrafo 1: popularização de assistentes de IA para geração de código.
% Parágrafo 2: problema -- código gerado é aceito com pouca revisão; complexidade alta
%   afeta manutenibilidade e testabilidade \cite{mccabe1976}.
% Parágrafo 3: lacuna -- pouca comparação controlada entre ferramentas usando métricas estruturais.
% Parágrafo 4: objetivo e questão de pesquisa (abaixo) + resumo da contribuição.
% Parágrafo 5: organização do artigo.

Este trabalho tem como objetivo comparar a complexidade estrutural do código gerado por
[IA A] e [IA B] para um mesmo conjunto de problemas de programação. Para isso, formulou-se
a seguinte questão de pesquisa:

\textbf{QP:} Há diferença na complexidade ciclomática das soluções geradas por [IA A] e
[IA B] para os mesmos problemas de programação?

Como questão secundária, investiga-se também o tamanho das soluções, medido em linhas de
código-fonte (SLOC).

%=====================================================================
\section{Fundamentação Teórica e Trabalhos Relacionados}
\subsection{Complexidade Ciclomática e LOC}
% Definição de McCabe: número de caminhos linearmente independentes, M = E - N + 2P.
% Na prática: 1 + número de pontos de decisão (if, for, while, and/or, except...).
% LOC/SLOC como medida de tamanho e suas limitações (dependente de estilo).
\cite{mccabe1976}

\subsection{Modelos de Linguagem para Geração de Código}
% Codex/HumanEval \cite{chen2021codex}; avaliação do Copilot com problemas do LeetCode
% \cite{nguyen2022copilot}. Acrescentar 2--3 trabalhos que vocês lerem (procurar no
% Google Scholar: "LLM generated code quality complexity").

%=====================================================================
\section{Metodologia}
O estudo segue as diretrizes de experimentação em Engenharia de Software de Wohlin et al.
\cite{wohlin2012} e caracteriza-se como um experimento \emph{in vitro} com
delineamento pareado: cada problema (unidade experimental) recebe os dois tratamentos.

\subsection{Hipóteses}
Sejam $\tilde{d}$ a mediana das diferenças pareadas de complexidade ciclomática total
($CC_A - CC_B$):
\begin{itemize}
  \item $H_0$: $\tilde{d} = 0$ (não há diferença na complexidade ciclomática entre as IAs);
  \item $H_1$: $\tilde{d} \neq 0$ (há diferença na complexidade ciclomática entre as IAs).
\end{itemize}
Hipóteses análogas foram definidas para SLOC ($H_{0,LOC}$ e $H_{1,LOC}$).

\subsection{Variáveis}
\begin{itemize}
  \item \textbf{Variável independente:} ferramenta de IA, com dois níveis: [IA A, modelo
  e versão] e [IA B, modelo e versão].
  \item \textbf{Variáveis dependentes:} (i) complexidade ciclomática total por solução
  (soma da complexidade do nível do módulo e de todas as funções e métodos);
  (ii) SLOC, linhas de código excluindo comentários e linhas em branco.
  \item \textbf{Variável de controle:} corretude da solução (aceita ou não pelo juiz online).
\end{itemize}

\subsection{Seleção dos Problemas}
% Explicar o critério ANTES de gerar os códigos. Sugestão:
Foram selecionados 30 problemas da plataforma [Beecrowd/LeetCode], estratificados por
dificuldade (10 fáceis, 10 médios e 10 difíceis), sorteados aleatoriamente dentro de cada
categoria. A lista completa, com identificadores, está no Apêndice/repositório [link].

\subsection{Instrumentação e Coleta}
Para cada problema, o enunciado foi enviado às duas IAs utilizando o mesmo
\emph{prompt} padronizado:
\begin{quote}\small
``Resolva o problema a seguir em Python 3. Responda apenas com o código, sem explicações.
[ENUNCIADO COMPLETO, incluindo formato de entrada e saída]''
\end{quote}
Para reduzir vieses: (i) cada problema foi enviado em uma nova conversa, sem histórico;
(ii) registrou-se a primeira resposta gerada, sem regenerações ou correções;
(iii) as versões dos modelos e as datas da coleta foram registradas
([DATA], [VERSÕES]); (iv) cada solução foi submetida ao juiz online para registrar a corretude.

As métricas foram extraídas automaticamente com a ferramenta Radon\footnote{\url{https://radon.readthedocs.io}}
(versão [X]), por meio de um script disponibilizado em [link do repositório].

\subsection{Análise dos Dados}
Foram calculadas estatísticas descritivas (média, desvio padrão, mediana, quartis, mínimo
e máximo) e gerados \emph{boxplots}. A normalidade das diferenças pareadas foi verificada
com o teste de Shapiro-Wilk. Por se tratar de amostras pareadas e de métricas discretas
com provável distribuição assimétrica, adotou-se o teste de postos sinalizados de Wilcoxon
(bicaudal, $\alpha = 0{,}05$). O tamanho de efeito foi medido pela correlação rank-bisserial
para amostras pareadas \cite{kerby2014} e pelo delta de Cliff, interpretado segundo os
limiares de Romano et al. \cite{romano2006}: desprezível ($|\delta| < 0{,}147$), pequeno ($< 0{,}33$),
médio ($< 0{,}474$) e grande ($\geq 0{,}474$). Como análise de sensibilidade, os testes
foram repetidos considerando apenas os pares em que ambas as soluções foram aceitas.

%=====================================================================
\section{Resultados}
\subsection{Corretude das Soluções}
% Taxa de aceitação de cada IA (saída do script). Ex.: "[IA A] teve X de 30 soluções aceitas..."

\subsection{Estatística Descritiva}
A Tabela~\ref{tab:descritiva} apresenta [...]. A Figura~\ref{fig:boxcc} [...].

\input{tabela_descritiva.tex}

\begin{figure}[ht]
\centering
\includegraphics[width=.48\textwidth]{boxplot_cc.pdf}
\includegraphics[width=.48\textwidth]{boxplot_loc.pdf}
\caption{Distribuição da complexidade ciclomática (esq.) e de SLOC (dir.) por IA}
\label{fig:boxcc}
\end{figure}

\subsection{Teste de Hipótese e Tamanho de Efeito}
% Copiar de resultados.txt: W, p, r, delta de Cliff e magnitude. Dizer se rejeita H0.
% Se rejeitar: indicar a DIREÇÃO (qual IA gera código mais complexo).
% Se não rejeitar: NÃO dizer que "as IAs são iguais", e sim que não houve evidência de diferença.

%=====================================================================
\section{Discussão}
% Responder a QP diretamente. Interpretar o tamanho de efeito na prática
% (ex.: diferença de 1 ponto de CC importa?). Relacionar com os trabalhos relacionados.
% Comentar a relação entre corretude e complexidade, se houver.

%=====================================================================
\section{Ameaças à Validade}
% Organizado segundo Wohlin et al. (2012).
\textbf{Validade interna.} Os modelos são não determinísticos: um mesmo \emph{prompt} pode
gerar soluções diferentes. Mitigou-se com protocolo fixo (nova conversa, primeira resposta),
mas a amostra representa uma única execução por problema. A redação do \emph{prompt}
também pode influenciar o estilo do código.

\textbf{Validade de construção.} Complexidade ciclomática e SLOC capturam apenas aspectos
estruturais da qualidade, não legibilidade, eficiência ou corretude. SLOC depende do estilo
de formatação. Soluções incorretas podem ser artificialmente simples ou complexas; por isso
realizou-se a análise de sensibilidade apenas com soluções aceitas.

\textbf{Validade externa.} Os resultados restringem-se a Python, a problemas de
programação competitiva (curtos e bem especificados) e às versões dos modelos na data da
coleta, não se generalizando para sistemas reais ou versões futuras. Os problemas podem
estar presentes nos dados de treinamento dos modelos (contaminação).

\textbf{Validade de conclusão.} Com $n = 30$ pares, o poder estatístico para detectar
efeitos pequenos é limitado. Métricas discretas geram empates, que são descartados pelo
teste de Wilcoxon; verificou-se a robustez com o método de Pratt. Testar duas métricas
aumenta a chance de erro do tipo I; SLOC é tratada como análise secundária.

%=====================================================================
\section{Aspectos Éticos}
O estudo não envolve participantes humanos nem dados pessoais. Os enunciados dos problemas
não são reproduzidos neste artigo, apenas seus identificadores, respeitando os termos de uso
das plataformas. Todos os dados brutos (códigos gerados, métricas e script) estão disponíveis
em [link], permitindo replicação. % Declarar aqui se usaram IA para ajudar na escrita.

%=====================================================================
\section{Conclusão}
% Resposta à QP em 1--2 frases, contribuição, limitações principais e trabalhos futuros
% (mais linguagens, múltiplas execuções por problema, métricas de manutenibilidade, outros modelos).

\bibliographystyle{sbc}
\bibliography{referencias}

\end{document}
